% © 2026 Ing. David Jaroš — CC BY-NC-ND 4.0
% UBT-AI-PROVENANCE-BEGIN
% schema: ubt-ai-provenance/v1
% tier: B_machine_verified
% ai_assistance: disclosed
% human_review: machine-verification
% editorial_responsibility: Ing. David Jaroš
% policy: ../../../../AI_PROVENANCE.md
% notice: Machine-verified against named sources or verifiers; individual attestation is not claimed.
% UBT-AI-PROVENANCE-END

\section{The central anticommutator metric}
\label{sec:canonical:biquaternion_metric}

Let $E_\mu=\mathcal N_0^{-1/2}D_\mu\Theta$ be the covariant UBT tetrad.  On the
classical Lorentz slice $W_L$, the fundamental quadratic identity is
\begin{equation}
 \boxed{E_\mu^\sharp E_\nu
       =g_{\mu\nu}\mathbf1+\Sigma_{\mu\nu},}
 \label{eq:canonical:full-product}
\end{equation}
where
\begin{align}
 g_{\mu\nu}\mathbf1
 &=\frac12(E_\mu^\sharp E_\nu+E_\nu^\sharp E_\mu),\\
 \Sigma_{\mu\nu}
 &=\frac12(E_\mu^\sharp E_\nu-E_\nu^\sharp E_\mu).
\end{align}
The first term is symmetric and central; the second is antisymmetric and
biquaternionic.  The classical metric is therefore defined intrinsically, not
by applying a trace, real-part operator, phase projector, or fiber average to a
noncentral object.

In components,
\begin{equation}
 g_{\mu\nu}=e_\mu{}^a e_\nu{}^b\eta_{ab},
 \qquad \eta_{ab}=\operatorname{diag}(-1,1,1,1).
\end{equation}
Nondegeneracy is equivalent to invertibility of $e_\mu{}^a$.  The local
Lorentzian signature follows immediately.

The full product in Eq.~\eqref{eq:canonical:full-product} remains relevant: its
antisymmetric part may carry spin/bivector information.  No part is discarded,
but no identification with a gauge field is claimed without a transformation
and dynamical derivation.
